\documentclass[11pt,a4paper]{article}
\usepackage{get-uebung}

\blatt{2}{Drahtwiderstand und Temperaturabhängigkeit}{Kap.~1.6--1.7}

\begin{document}

\blattkopf{
  \item den Widerstand eines Leiters aus Länge, Querschnitt und Material mit $R = \rho\, l/A = l/(\kappa A)$ berechnen,
  \item einen Leiterquerschnitt nach einem zulässigen Spannungsfall auswählen,
  \item die lineare Näherung $R(\vartheta) \approx R_{20}\,(1 + \alpha_{20}\,\Delta\vartheta)$ anwenden und nach jeder Größe umstellen,
  \item einen Temperaturkoeffizienten aus Messwerten bestimmen und erklären, warum er von der Bezugstemperatur abhängt.
}

% -------------------------------------------------------------------
\begin{aufgabe}{Das Verlängerungskabel}{15}
Ein Wasserkocher ($\SI{230}{V}$, $\SI{2000}{W}$) wird über eine $\SI{25}{m}$ lange Kabeltrommel betrieben. Das Kabel hat zwei Kupferadern (Hin- und Rückleiter) mit je $A = \SI{1.5}{mm^2}$ Querschnitt; $\kappa_\text{Cu} = \SI{56}{m/(\ohm.mm^2)}$.
\begin{teilaufgaben}
  \item Berechnen Sie den Widerstand $R_\text{L}$ und den Leitwert $G_\text{L}$ der gesamten Leitung.
  \item Wie groß sind näherungsweise der Spannungsfall $\Delta U$ auf der Leitung (absolut und in Prozent von $\SI{230}{V}$) und die Verlustleistung im Kabel? Rechnen Sie dazu mit dem Nennstrom des Wasserkochers.
  \item Der Spannungsfall soll höchstens $\SI{1}{\percent}$ betragen. Welcher Querschnitt ist mindestens nötig? Wählen Sie aus den Normquerschnitten $\SI{1.5}{mm^2}$, $\SI{2.5}{mm^2}$, $\SI{4}{mm^2}$, $\SI{6}{mm^2}$.
\end{teilaufgaben}
\begin{tipp}Der Strom muss zum Gerät hin \emph{und} wieder zurück fließen. Wie lang ist also der Leiter, den der Strom durchläuft?\end{tipp}
\end{aufgabe}

\begin{loesung}
\schritt{a) Widerstand und Leitwert}
\leitfrage{Welche Länge $l$ gehört in die Formel -- \SI{25}{m} oder \SI{50}{m}?}
Der Strom durchläuft Hin- und Rückleiter, also $l = 2\cdot\SI{25}{m} = \SI{50}{m}$. Mit $R = l/(\kappa A)$:
\[
  R_\text{L} = \frac{\SI{50}{m}}{\SI{56}{\frac{m}{\ohm\,mm^2}}\cdot\SI{1.5}{mm^2}} \approx \ergebnis{\SI{0.595}{\ohm}},
  \qquad G_\text{L} = \frac{1}{R_\text{L}} = \ergebnis{\SI{1.68}{S}}
\]
Die Einheiten kürzen sich: $\frac{\mathrm{m}}{\frac{\mathrm{m}}{\Omega\,\mathrm{mm^2}}\cdot\mathrm{mm^2}} = \Omega$. Deshalb ist die Angabe von $\kappa$ in $\si{m/(\ohm.mm^2)}$ so praktisch: $l$ in Metern und $A$ in Quadratmillimetern direkt einsetzen.
\schritt{b) Spannungsfall und Verluste}
Nennstrom: $I = P/U = \SI{2000}{W}/\SI{230}{V} \approx \SI{8.70}{A}$ (der kleine Leitungswiderstand ändert ihn kaum).
\[
  \Delta U = R_\text{L}\cdot I \approx \SI{0.595}{\ohm}\cdot\SI{8.70}{A} \approx \ergebnis{\SI{5.18}{V}} \;\hat{=}\; \SI{2.25}{\percent},
  \qquad P_\text{V} = I^2 R_\text{L} \approx \ergebnis{\SI{45}{W}}
\]
Das Kabel heizt also mit $\SI{45}{W}$ -- auf der aufgerollten Trommel kann das gefährlich warm werden!
\schritt{c) Querschnitt auswählen}
\leitfrage{Stellen Sie $\Delta U = I\cdot l/(\kappa A)$ nach $A$ um.}
Zulässig sind $\Delta U_\text{max} = 0{,}01\cdot\SI{230}{V} = \SI{2.3}{V}$:
\[
  A_\text{min} = \frac{l\cdot I}{\kappa\cdot\Delta U_\text{max}} = \frac{\SI{50}{m}\cdot\SI{8.70}{A}}{\SI{56}{\frac{m}{\ohm\,mm^2}}\cdot\SI{2.3}{V}} \approx \SI{3.38}{mm^2}
  \quad\Rightarrow\quad \ergebnis{A = \SI{4}{mm^2}}
\]
Es muss der nächst\emph{größere} Normquerschnitt gewählt werden, sonst wird die Grenze überschritten.
\typfehler{Nur die einfache Kabellänge einsetzen (Ergebnis halb so groß). Den Querschnitt auf den \emph{nächstgelegenen} statt auf den nächstgrößeren Normwert runden. $\mathrm{mm^2}$ in $\mathrm{m^2}$ umrechnen, obwohl $\kappa$ schon in $\si{m/(\ohm.mm^2)}$ gegeben ist -- dann muss man konsequent \emph{alles} umrechnen ($\SI{1}{mm^2} = \SI{e-6}{m^2}$).}
\end{loesung}

% -------------------------------------------------------------------
\begin{aufgabe}{Kupfer oder Aluminium?}{12}
Für eine Freileitung soll ein Leiter aus Kupfer durch einen Leiter aus Aluminium \emph{mit gleichem Widerstand und gleicher Länge} ersetzt werden.

\gegeben{$\rho_\text{Cu} = \SI{0.0178}{\ohm.mm^2/m}$, $\rho_\text{Al} = \SI{0.0286}{\ohm.mm^2/m}$, Dichten $\varrho_{m,\text{Cu}} = \SI{8.96}{g/cm^3}$, $\varrho_{m,\text{Al}} = \SI{2.70}{g/cm^3}$.}
\begin{teilaufgaben}
  \item Um welchen Faktor muss der Querschnitt des Aluminiumleiters größer sein? Um welchen Faktor wächst damit der Durchmesser eines runden Leiters?
  \item Wie verhält sich die Masse des Aluminiumleiters zu der des Kupferleiters?
  \item Begründen Sie damit, warum Hochspannungsfreileitungen meist aus Aluminium (mit Stahlseele) bestehen, Kabel in Gebäuden dagegen aus Kupfer.
\end{teilaufgaben}
\end{aufgabe}

\begin{loesung}
\schritt{a) Querschnitt und Durchmesser}
\leitfrage{Setzen Sie $R_\text{Al} = R_\text{Cu}$ mit $R = \rho\,l/A$ an. Was kürzt sich heraus?}
\[
  \rho_\text{Al}\frac{l}{A_\text{Al}} = \rho_\text{Cu}\frac{l}{A_\text{Cu}}
  \quad\Rightarrow\quad
  \frac{A_\text{Al}}{A_\text{Cu}} = \frac{\rho_\text{Al}}{\rho_\text{Cu}} = \frac{0{,}0286}{0{,}0178} \approx \ergebnis{1{,}61}
\]
Wegen $A = \pi d^2/4$ ist $A \sim d^2$, also
$\dfrac{d_\text{Al}}{d_\text{Cu}} = \sqrt{1{,}61} \approx \ergebnis{1{,}27}$.
\schritt{b) Massenverhältnis}
Masse $=$ Dichte $\cdot$ Volumen $= \varrho_m\cdot A\cdot l$. Bei gleicher Länge:
\[
  \frac{m_\text{Al}}{m_\text{Cu}} = \frac{\varrho_{m,\text{Al}}}{\varrho_{m,\text{Cu}}}\cdot\frac{A_\text{Al}}{A_\text{Cu}}
  = \frac{2{,}70}{8{,}96}\cdot 1{,}61 \approx \ergebnis{0{,}48}
\]
Der Aluminiumleiter ist trotz größeren Querschnitts nur etwa \emph{halb so schwer}.
\schritt{c) Bewertung}
Bei Freileitungen zählt das Gewicht (Masten, Seilspannung) und der Preis; der größere Durchmesser stört in der Luft nicht. In Gebäuden ist der Platz in Rohren und Dosen knapp, und Kupfer lässt sich zuverlässiger klemmen -- deshalb Kupfer.
\typfehler{Den Kehrwert bilden ($A_\text{Al}/A_\text{Cu} = \rho_\text{Cu}/\rho_\text{Al}$): Das schlechter leitende Material braucht den \emph{größeren} Querschnitt. Faktor für $A$ und für $d$ verwechseln.}
\end{loesung}

% -------------------------------------------------------------------
\begin{aufgabe}{Wann verdoppelt sich ein Widerstand?}{13}
Die Wicklung eines Elektromotors besteht aus Kupferdraht ($\alpha_{20} = \SI{3.93e-3}{1/K}$). Verwenden Sie die lineare Näherung
$R(\vartheta) \approx R_{20}\,\bigl(1 + \alpha_{20}\,(\vartheta - \SI{20}{\celsius})\bigr)$.
\begin{teilaufgaben}
  \item Bei welcher Temperatur hätte die Wicklung (nach der linearen Näherung) den doppelten Widerstand wie bei $\SI{20}{\celsius}$? Hängt das Ergebnis vom Wert $R_{20}$ ab?
  \item Die Wicklungstemperatur eines Motors lässt sich nicht direkt messen, weil die Wicklung im Inneren liegt. Man misst deshalb den Widerstand: Im kalten Zustand ($\SI{20}{\celsius}$) beträgt er $\SI{2.00}{\ohm}$, direkt nach einer Stunde Betrieb $\SI{2.40}{\ohm}$. Wie warm ist die Wicklung?
  \item Welchen Strom nimmt der Motor beim Anlauf an $\SI{24}{V}$ im kalten und im warmen Zustand auf, wenn man (im Stillstand) nur den Wicklungswiderstand betrachtet?
\end{teilaufgaben}
\end{aufgabe}

\begin{loesung}
\schritt{a) Verdopplungstemperatur}
\leitfrage{Setzen Sie $R(\vartheta) = 2R_{20}$ in die Näherung ein. Was passiert mit $R_{20}$?}
\[
  2R_{20} = R_{20}\,(1 + \alpha_{20}\,\Delta\vartheta)
  \quad\Rightarrow\quad 1 = \alpha_{20}\,\Delta\vartheta
  \quad\Rightarrow\quad \Delta\vartheta = \frac{1}{\alpha_{20}} = \frac{1}{\SI{3.93e-3}{1/K}} \approx \SI{254.5}{K}
\]
\[
  \vartheta = \SI{20}{\celsius} + \SI{254.5}{K} \approx \ergebnis{\SI{274.5}{\celsius}}
\]
$R_{20}$ kürzt sich heraus: Das Ergebnis gilt für \emph{jeden} Kupferleiter. (In Wahrheit ist $R(\vartheta)$ nicht exakt linear, so weit von $\SI{20}{\celsius}$ entfernt ist die Näherung nur grob.)
\schritt{b) Wicklungstemperatur aus Widerstandsmessung}
Umstellen nach $\Delta\vartheta$:
\[
  \Delta\vartheta = \frac{1}{\alpha_{20}}\left(\frac{R(\vartheta)}{R_{20}} - 1\right)
  = \frac{1}{\SI{3.93e-3}{1/K}}\left(\frac{\SI{2.40}{\ohm}}{\SI{2.00}{\ohm}} - 1\right)
  = \frac{0{,}2}{\SI{3.93e-3}{1/K}} \approx \SI{50.9}{K}
\]
\[
  \vartheta \approx \SI{20}{\celsius} + \SI{50.9}{K} \approx \ergebnis{\SI{70.9}{\celsius}}
\]
\schritt{c) Anlaufströme}
\[
  I_\text{kalt} = \frac{\SI{24}{V}}{\SI{2.00}{\ohm}} = \ergebnis{\SI{12}{A}}, \qquad
  I_\text{warm} = \frac{\SI{24}{V}}{\SI{2.40}{\ohm}} = \ergebnis{\SI{10}{A}}
\]
Der kalte Motor zieht beim Anlauf also $\SI{20}{\percent}$ mehr Strom.
\typfehler{$\vartheta$ statt $\Delta\vartheta$ einsetzen, also das Ergebnis $\SI{254.5}{\celsius}$ angeben -- die Näherung arbeitet mit der \emph{Differenz} zur Bezugstemperatur. Temperaturdifferenzen werden in $\si{K}$ angegeben, absolute Temperaturen hier in $\si{\celsius}$.}
\end{loesung}

% -------------------------------------------------------------------
\ifloesung\else\newpage\fi
\begin{aufgabe}{Temperaturkoeffizient aus Messwerten}{17}
Ein Platin-Temperatursensor wird in einem Wasserbad vermessen:
\begin{center}
\begin{tabular}{l*{5}{c}}
\toprule
$\vartheta$ in $\si{\celsius}$ & 0 & 20 & 40 & 60 & 80\\
$R$ in $\si{\ohm}$ & \num{100.00} & \num{107.79} & \num{115.54} & \num{123.24} & \num{130.90}\\
\bottomrule
\end{tabular}
\end{center}
\begin{teilaufgaben}
  \item Tragen Sie die Messwerte in ein Diagramm $R(\vartheta)$ ein. Ist der Zusammenhang exakt linear?
  \item Bestimmen Sie die Steigung $\mathrm{d}R/\mathrm{d}\vartheta$ bei $\SI{20}{\celsius}$ näherungsweise aus den beiden benachbarten Messpunkten und daraus den Temperaturkoeffizienten $\alpha_{20}$.
  \item Berechnen Sie mit der linearen Näherung den Widerstand bei $\SI{80}{\celsius}$ und vergleichen Sie mit dem Messwert. Wie groß ist der Fehler in Ohm und in Prozent?
  \item Dieselbe Gerade soll mit der Bezugstemperatur $\SI{0}{\celsius}$ geschrieben werden: $R(\vartheta) \approx R_0\,(1 + \alpha_0\,\vartheta)$. Zeigen Sie, dass $\alpha_0 = \dfrac{\alpha_{20}}{1 - \alpha_{20}\cdot\SI{20}{K}}$ gilt, und berechnen Sie $\alpha_0$. Warum ist $\alpha_0 > \alpha_{20}$?
\end{teilaufgaben}
\begin{tipp}Laut Skript ist $\alpha_{\vartheta_0} = \dfrac{1}{R(\vartheta_0)}\cdot\dfrac{\mathrm{d}R}{\mathrm{d}\vartheta}\Big|_{\vartheta_0}$ -- die Steigung wird auf den Widerstand bei der Bezugstemperatur bezogen.\end{tipp}
\end{aufgabe}

\begin{loesung}
\schritt{a) Diagramm}
\begin{center}
\begin{tikzpicture}
\begin{axis}[width=8cm,height=5.2cm,xlabel={$\vartheta$ in $\si{\celsius}$},ylabel={$R$ in $\si{\ohm}$},
  xmin=-5,xmax=85,ymin=95,ymax=135,grid=major,legend pos=north west,legend style={font=\small}]
  \addplot[only marks,mark=*,hsublau] coordinates {(0,100.00) (20,107.79) (40,115.54) (60,123.24) (80,130.90)};
  \addlegendentry{Messwerte}
  \addplot[loesungfarbe,thick,domain=-5:85] {107.79*(1+0.0036042*(x-20))};
  \addlegendentry{Tangente bei $\SI{20}{\celsius}$}
\end{axis}
\end{tikzpicture}
\end{center}
Die Differenzen zwischen benachbarten Werten sind $7{,}79$; $7{,}75$; $7{,}70$; $\SI{7.66}{\ohm}$ -- sie werden leicht kleiner. Die Kennlinie ist also \emph{fast}, aber nicht exakt linear (leicht nach unten gekrümmt).
\schritt{b) Steigung und $\alpha_{20}$}
\leitfrage{Welche zwei Messpunkte liegen symmetrisch um $\SI{20}{\celsius}$?}
Die Punkte bei $\SI{0}{\celsius}$ und $\SI{40}{\celsius}$:
\[
  \frac{\mathrm{d}R}{\mathrm{d}\vartheta}\Big|_{20} \approx \frac{\SI{115.54}{\ohm} - \SI{100.00}{\ohm}}{\SI{40}{K}} = \SI{0.3885}{\ohm/K}
\]
\[
  \alpha_{20} = \frac{1}{R_{20}}\cdot\frac{\mathrm{d}R}{\mathrm{d}\vartheta}\Big|_{20} = \frac{\SI{0.3885}{\ohm/K}}{\SI{107.79}{\ohm}} \approx \ergebnis{\SI{3.60e-3}{1/K}}
\]
\schritt{c) Vorhersage bei $\SI{80}{\celsius}$}
\[
  R(\SI{80}{\celsius}) \approx \SI{107.79}{\ohm}\cdot\bigl(1 + \SI{3.604e-3}{1/K}\cdot\SI{60}{K}\bigr) \approx \ergebnis{\SI{131.10}{\ohm}}
\]
Fehler: $\SI{131.10}{\ohm} - \SI{130.90}{\ohm} = \SI{0.20}{\ohm}$, also etwa $\SI{0.15}{\percent}$. Je weiter man sich von der Bezugstemperatur entfernt, desto größer wird der Fehler (Abb.~1.4 im Skript).
\schritt{d) Umrechnung auf $\SI{0}{\celsius}$}
\leitfrage{Beide Darstellungen beschreiben dieselbe Gerade. Welchen Wert $R_0$ liefert die $\SI{20}{\celsius}$-Gerade bei $\vartheta = \SI{0}{\celsius}$, und welche Steigung hat sie?}
Gerade bei $\SI{0}{\celsius}$: $R_0 = R_{20}\,(1 - \alpha_{20}\cdot\SI{20}{K})$. Die Steigung ist überall gleich: $R_{20}\,\alpha_{20} = R_0\,\alpha_0$. Also
\[
  \alpha_0 = \frac{R_{20}\,\alpha_{20}}{R_0} = \frac{\alpha_{20}}{1 - \alpha_{20}\cdot\SI{20}{K}}
  = \frac{\SI{3.604e-3}{1/K}}{1 - 0{,}0721} \approx \ergebnis{\SI{3.88e-3}{1/K}}
\]
$\alpha_0$ ist größer, weil dieselbe Steigung auf den \emph{kleineren} Widerstand $R_0 < R_{20}$ bezogen wird.
\typfehler{Die Steigung nicht auf $R_{20}$ beziehen (dann hat $\alpha$ die falsche Einheit $\si{\ohm/K}$). Glauben, $\alpha$ sei eine reine Materialkonstante unabhängig von der Bezugstemperatur -- deshalb steht der Index immer dabei.}
\end{loesung}

\end{document}
